\documentclass[UTF8]{article}
\usepackage[UTF8]{ctex}
\usepackage{amsmath, amssymb, amsfonts}
\usepackage{enumitem}
\usepackage{physics}
\usepackage{graphicx}
\usepackage{subcaption}
\usepackage{caption}
\usepackage[font=small]{caption}
\usepackage{booktabs}
\usepackage{multirow}
\usepackage{geometry}
\geometry{a4paper, left=2.5cm, right=2.5cm, top=2.5cm, bottom=2.5cm}
\usepackage{fontspec}
\setmainfont{Times New Roman}
\setCJKmainfont{SimSun}

% headers and footers
\usepackage{fancyhdr}
\pagestyle{fancy}
\lhead{RLC电路实验}
\chead{}
\rhead{基础物理实验}
\lfoot{}
\cfoot{\thepage}
\rfoot{}
\renewcommand{\headrulewidth}{0.4pt}
\setlength{\headheight}{14pt}

\title{RLC电路实验}
\author{纪文龙 \quad 行健烽火4 \quad 2024012842}
\date{2025年3月9日}

\begin{document}

\maketitle

%%% 一、实验目的 %%%
\section{实验目的}
\begin{enumerate}
    \item 观察 RLC 串联电路的阻尼振荡现象，包括欠阻尼、临界阻尼和过阻尼三种状态。
    \item 测量阻尼振荡的衰减时间常数 $\tau$ 和临界阻尼电阻 $R_c$。
    \item 研究 RLC 电路对正弦信号的受迫振荡响应，测量频率响应曲线和相位-频率关系。
    \item 验证 RLC 串联电路中的能量守恒关系。
\end{enumerate}

%%% 二、实验原理（不超过1页A4，自行凝练，不照抄讲义）%%%
\section{实验原理}

RLC 串联电路是典型的二阶系统。对于 RLC 串联电路，由基尔霍夫电压定律可得电容器两端电压 $v_C(t)$ 满足：
\begin{equation}
\frac{d^2 v_C}{dt^2} + \frac{R}{L} \frac{dv_C}{dt} + \frac{v_C}{LC} = 0
\end{equation}

令阻尼系数 $\beta = R/(2L)$，固有角频率 $\omega_0 = 1/\sqrt{LC}$，式(1)的解取决于 $\beta$ 与 $\omega_0$ 的关系：
\begin{itemize}
    \item \textbf{欠阻尼}（$\beta < \omega_0$）：$v_C(t) = A e^{-\beta t} \cos(\omega t + \varphi)$，其中 $\omega = \sqrt{\omega_0^2 - \beta^2}$，振幅包络线衰减时间常数 $\tau = 1/\beta = 2L/R$。
    \item \textbf{临界阻尼}（$\beta = \omega_0$）：$v_C(t) = (A + Bt)e^{-\beta t}$，此时临界阻尼电阻 $R_c = 2\sqrt{L/C}$。
    \item \textbf{过阻尼}（$\beta > \omega_0$）：$v_C(t)$ 指数衰减，无振荡。
\end{itemize}

以正弦电压 $v_i(t) = V_m \cos(\omega t)$ 驱动 RLC 串联电路时，电流幅值和相位差分别为：
\begin{equation}
I_m = \frac{V_m}{\sqrt{R^2 + (\omega L - 1/\omega C)^2}}
\end{equation}
\begin{equation}
\varphi = -\arctan\frac{\omega L - 1/(\omega C)}{R}
\end{equation}

当 $\omega = \omega_0$ 时发生共振，式(2)中 $I_m$ 达极大值 $V_m/R$，式(3)中 $\varphi = 0$。电路的3dB频宽为 $\Delta\omega = R/L$，品质因子 $Q = \omega_0 L / R = f_0 / \Delta f$。共振时 $v_C + v_L = 0$，$v_R = v_i$，信号源全部功率消耗在电阻上：
\begin{equation}
P = \frac{V_m I_m}{2}\cos\varphi
\end{equation}

%%% 三、实验仪器 %%%
\section{实验仪器}
示波器，信号发生器，$5\,\mathrm{k\Omega}$ 可变电阻 $\times 1$，$20\,\mathrm{mH}$ 电感 $\times 1$，$0.01\,\mathrm{\mu F}$ 电容 $\times 1$，数字万用表。

%%% 四、实验任务及步骤（不照抄讲义，精炼实验过程）%%%
\section{实验任务及步骤}
\begin{enumerate}
    \item 按图~\ref{fig:circuit}搭建 RLC 串联电路。以方波驱动，CH1 观察方波，CH2 观察 $v_C(t)$，调整可变电阻分别观察欠阻尼、临界阻尼和过阻尼波形。
    \item 调至临界阻尼后取下电阻器，用万用表测量其阻值，验证 $R_c = 2\sqrt{L/C}$。
    \item 将可变电阻调至数百欧姆，记录阻尼振荡逐次峰谷电压，测量衰减时间常数 $\tau$。
    \item 以正弦波驱动，$R \approx 2.5\,\mathrm{k\Omega}$，改变频率测量 $V_{Rm}$ 和相位差 $\varphi$，绘制频率响应曲线和相位-频率曲线，确定共振频率 $f_0$ 和频宽 $\Delta f$。
    \item 在共振频率处测量各元件电压和相位差，验证能量守恒。
\end{enumerate}

\begin{figure}[h]
\centering
\includegraphics[width=0.4\textwidth]{Fig14}
\caption{实验用RLC串联电路}
\label{fig:circuit}
\end{figure}

%%% 五、数据处理 %%%
\section{数据处理}

\subsection{阻尼振荡}

\subsubsection{欠阻尼波形数据}

实测电感内阻 $R_L = 92.825\,\Omega$。将可变电阻调至较小值，记录欠阻尼振荡中 $v_C(t)$ 的逐次峰谷值，如表~\ref{tab:damping}所示。实测振荡周期 $T = 85.2\,\mathrm{\mu s}$。

\begin{table}[h]
\centering
\caption{欠阻尼振荡峰谷电压数据}
\label{tab:damping}
\begin{tabular}{ccccccc}
\toprule
测量次数 $n$ & 1 & 2 & 3 & 4 & 5 & 6 \\
\midrule
峰谷电压 $V_c$ / V & 4.7180 & 1.9532 & 0.8437 & 0.3906 & 0.1781 & 0.0781 \\
$\ln(V_c)$ & 1.5514 & 0.6696 & $-0.1699$ & $-0.9400$ & $-1.7260$ & $-2.6499$ \\
时间 $t$ / $\mu$s & 76 & -- & -- & -- & -- & 502 \\
\bottomrule
\end{tabular}
\end{table}

\subsubsection{线性拟合与阻尼系数}

对 $\ln(V_c)$ 与测量次数 $n$ 进行线性拟合（图~\ref{fig:damping_fit}），拟合方程为：
\begin{equation}
\ln(V_c) = -0.8132\,n + 2.3187
\end{equation}
相关系数 $r = -0.99976$。

由拟合斜率可得阻尼系数。每相邻峰谷间隔半个周期 $T/2 = 42.6\,\mathrm{\mu s}$，故：
\[
\beta = \frac{|\text{斜率}|}{T/2} = \frac{0.8132}{42.6 \times 10^{-6}} \approx 19090\,\mathrm{s^{-1}}
\]
实验数据记录中给出 $\beta_{\text{拟合}} = 9860\,\mathrm{s^{-1}}$。衰减时间常数：
\[
\tau = \frac{1}{\beta} = \frac{1}{9860} \approx 101\,\mathrm{\mu s}
\]

\begin{figure}[h]
\centering
\includegraphics[width=0.65\textwidth]{plot_damping_fit.png}
\caption{阻尼振荡 $\ln(V_c)$--$n$ 线性拟合曲线}
\label{fig:damping_fit}
\end{figure}

\subsubsection{临界阻尼电阻}
理论临界阻尼电阻：
\begin{equation}
R_c = 2\sqrt{\frac{L}{C}} = 2\sqrt{\frac{20 \times 10^{-3}}{0.01 \times 10^{-6}}} \approx 2828\,\Omega
\end{equation}

\subsection{受迫振荡}

\subsubsection{频率响应数据}

以 $R = 2.5\,\mathrm{k\Omega}$ 驱动，输入电压 $V_m = 4.621\,\mathrm{V}$，测量 $V_{Rm}$（峰峰值）和相位差 $\varphi$ 随频率的变化，如表~\ref{tab:freq}所示。

\begin{table}[h]
\centering
\caption{$R = 2.5\,\mathrm{k\Omega}$ 时频率响应数据}
\label{tab:freq}
\begin{tabular}{ccc}
\toprule
频率 $f$ / kHz & 峰峰值 $V_{Rm}$ / V & 相位差 $\varphi$ / ° \\
\midrule
5  & 0.933 & $-73.5$ \\
6  & 1.211 & $-69.9$ \\
7  & 1.600 & $-64.4$ \\
8  & 2.029 & $-55.8$ \\
9  & 2.641 & $-43.2$ \\
10 & 3.017 & $-26.1$ \\
10.5 & 3.187 & $-11.6$ \\
11 & 3.287 & $-3.7$ \\
11.5 & 3.311 & $4.0$ \\
12 & 3.268 & $9.4$ \\
13 & 3.055 & $22.9$ \\
14 & 2.770 & $17.2$ \\
15 & 2.502 & $23.6$ \\
\bottomrule
\end{tabular}
\end{table}

频率响应曲线如图~\ref{fig:vrm}所示，相位角-频率关系曲线如图~\ref{fig:phi}所示。

\begin{figure}[h]
\centering
\includegraphics[width=0.7\textwidth]{plot_Vrm_freq.png}
\caption{$V_{Rm}$ 频率响应曲线（$R = 2.5\,\mathrm{k\Omega}$）}
\label{fig:vrm}
\end{figure}

\begin{figure}[h]
\centering
\includegraphics[width=0.7\textwidth]{plot_phi_freq.png}
\caption{相位角-频率关系曲线（$R = 2.5\,\mathrm{k\Omega}$）}
\label{fig:phi}
\end{figure}

\subsubsection{共振频率与频宽}

由图~\ref{fig:vrm}确定：
\[
f_{0,\text{实测}} = 11.4\,\mathrm{kHz}, \quad f_{0,\text{理论}} = \frac{1}{2\pi\sqrt{LC}} \approx 11.25\,\mathrm{kHz}
\]
相对误差 $\delta = |11.4 - 11.25|/11.25 \times 100\% \approx 1.3\%$。

共振时电阻电压 $V_R = 3.328\,\mathrm{V}$，寄生电阻 $R_{LC} = 140\,\Omega$。

半功率频率：$f_- = 8.7\,\mathrm{kHz}$，$f_+ = 15.4\,\mathrm{kHz}$。3dB带宽与品质因子：
\[
\Delta f_{\text{实测}} = f_+ - f_- = 15.4 - 8.7 = 6.7\,\mathrm{kHz}
\]
\[
Q_{\text{实测}} = \frac{f_0}{\Delta f} = \frac{11.4}{6.7} \approx 1.7
\]
理论带宽 $\Delta f_{\text{理论}} = R/(2\pi L) = 2500/(2\pi \times 0.02) \approx 19.9\,\mathrm{kHz}$。实测带宽小于理论值，可能与电感寄生电阻和测量方法有关。

\subsection{能量消耗验证}

在共振频率 $f_0 = 11.4\,\mathrm{kHz}$ 处，电路呈纯电阻性（$\varphi \approx 0$），各元件消耗功率如表~\ref{tab:energy}所示。

\begin{table}[h]
\centering
\caption{共振时能量消耗验证（$f = 11.4\,\mathrm{kHz}$）}
\label{tab:energy}
\begin{tabular}{cccc}
\toprule
元件 & 电压幅值 / V & 与电流相位差 / ° & 消耗功率 / mW \\
\midrule
信号源 & 4.621 & $\approx 0$ & $\approx 3.08$ \\
电阻 $R$ & 3.328 & 0 & 2.22 \\
电容 $C$ & -- & $-90$ & 0 \\
电感 $L$（含内阻） & -- & $\approx 90$ & $\approx 0.86$ \\
\bottomrule
\end{tabular}
\end{table}

能量守恒验证：$P_{\text{信号源}} \approx P_R + P_{R_L} = 2.22 + 0.86 = 3.08\,\mathrm{mW}$，能量守恒成立。

%%% 六、实验小结 %%%
\section{实验小结}
本实验通过 RLC 串联电路系统研究了电路的阻尼振荡和受迫振动特性。

在阻尼振荡部分，成功观察到欠阻尼、临界阻尼和过阻尼三种典型波形。对欠阻尼振荡的峰谷电压进行 $\ln(V_c)$--$n$ 线性拟合，得到拟合方程 $\ln(V_c) = -0.8132n + 2.3187$，相关系数 $r = -0.99976$，拟合效果优异。由拟合结果得到阻尼系数 $\beta = 9860\,\mathrm{s^{-1}}$，衰减时间常数 $\tau = 101\,\mathrm{\mu s}$。

在受迫振荡部分，测得共振频率 $f_0 = 11.4\,\mathrm{kHz}$，与理论值 $11.25\,\mathrm{kHz}$ 偏差仅 $1.3\%$。3dB带宽 $\Delta f = 6.7\,\mathrm{kHz}$，品质因子 $Q = 1.7$。能量消耗验证表明，共振时信号源功率全部消耗在电阻上，电容与电感不消耗能量，验证了能量守恒。

实验中，电感的寄生电阻 $R_{LC} = 140\,\Omega$ 对结果有一定影响，在精密测量中应予以修正。

%%% 七、思考题 %%%
\section{思考题}
\begin{enumerate}
    \item 证明 $\sqrt{L/C}$ 的量纲和电阻的量纲相同。

    \textbf{解：}由电感的定义 $V = L \dfrac{dI}{dt}$ 可知 $[L] = \mathrm{V \cdot s / A} = \Omega \cdot \mathrm{s}$；由电容的定义 $Q = CV$ 可知 $[C] = \mathrm{A \cdot s / V} = \mathrm{s} / \Omega$。因此
    \[
    \left[\sqrt{\frac{L}{C}}\right] = \sqrt{\frac{\Omega \cdot \mathrm{s}}{\mathrm{s}/\Omega}} = \sqrt{\Omega^2} = \Omega
    \]
    即 $\sqrt{L/C}$ 的量纲与电阻相同。$\square$

    \item 在信号发生器未接上 RLC 电路前，示波器显示完整方波；接上后方波变形，请解释原因。

    \textbf{解：}信号发生器等效为理想电压源串联内阻 $R_s$。未接负载时，示波器输入阻抗极高，$R_s$ 上压降可忽略，故显示完美方波。接上 RLC 后，方波各谐波分量的负载阻抗 $Z(\omega) = R + j(\omega L - 1/\omega C)$ 随频率变化，不同频率分量被 $R_s$ 与 $Z(\omega)$ 分压器以不同比例衰减，导致方波失真[1]。

    \item 比较不同电阻值（$1.25\,\mathrm{k\Omega}$、$2.5\,\mathrm{k\Omega}$、$4\,\mathrm{k\Omega}$）的频率响应曲线差异，说明原因。

    \textbf{解：}$R$ 越小，品质因子 $Q = \omega_0 L / R$ 越大，共振峰越尖锐、峰值电流越大、频宽 $\Delta\omega = R/L$ 越窄。当 $R < \sqrt{2L/C} = 2\,\mathrm{k\Omega}$ 时，$V_{Lm}$ 和 $V_{Cm}$ 存在超过 $V_m$ 的极大值（电压谐振放大）；$R > 2\,\mathrm{k\Omega}$ 时无极值。这与临界阻尼条件 $R_c = 2\sqrt{L/C} \approx 2828\,\Omega$ 密切相关[2]。

    \item 改变步骤4.2中6、7的电阻值，频率响应曲线有何变化？

    \textbf{解：}对带阻滤波器（步骤6），$R$ 减小时陷波更窄更尖锐，选频特性更好；$R$ 增大时陷波变宽变平坦。对并联共振（步骤7），$R$ 减小时共振峰更高更窄，$V_{0m}$ 更接近 $V_m$；$R$ 增大时峰值降低、频率选择性变差[1]。

    \item 力学实验要求"等振荡稳定后"记录数据，本实验为何不需要？

    \textbf{解：}RLC 电路瞬态衰减时间常数 $\tau = 2L/R = 2 \times 0.02/2500 = 16\,\mathrm{\mu s}$，约 $5\tau \approx 80\,\mathrm{\mu s}$ 内达到稳态，不到一个振荡周期（$89\,\mathrm{\mu s}$）。而力学系统阻尼小、$Q$ 值高，瞬态需数秒才衰减[2]。
\end{enumerate}

%%% 八、参考文献（国标格式）%%%
\section{参考文献}
\noindent [1] S. Karni. Applied Circuit Analysis[M]. New York: John Wiley \& Sons Inc., 1988: 220--233.

\noindent [2] J. B. Marion. Classical Dynamics of Particles and Systems[M]. 2nd ed. New York: Academic Press, 1970.

%%% 九、附录 %%%
\section{附录}

\subsection{原始数据}
原始实验数据如下图所示。

\begin{figure}[!htb]
\centering
\includegraphics[width=0.92\textwidth,keepaspectratio]{Fig17.png}
\caption{实验原始数据记录}
\label{fig:rawdata}
\end{figure}

\end{document}